\section{Methods}
\label{sec:methods}

We organize our analysis around four complementary perspectives: statistical testing, sentiment analysis, network analysis, and predictive modeling.

\subsection{Statistical Analysis}

For each interaction feature, we compare the distributions between biased and non-biased articles using the Mann-Whitney U test (non-parametric, appropriate given non-normal distributions).
Effect sizes are reported as Cohen's $d$ and rank-biserial correlation $r$.
Correlations between bias probability and interaction features are computed using both Pearson (linear) and Spearman (monotonic) coefficients.

We extend the basic comparison with advanced karma distribution features.
For each article's comment karma sequence $\{k_1, \ldots, k_n\}$, we compute:
\begin{itemize}
    \item \textbf{Entropy}: $H = -\sum_i p_i \log_2 p_i$, where $p_i$ is the probability of binned karma value $i$. Higher entropy indicates more dispersed (controversial) karma distributions.
    \item \textbf{Gini coefficient}: Measures inequality in absolute karma values. Higher Gini indicates that a few comments dominate the karma distribution.
    \item \textbf{Bimodality coefficient}: $BC = (\gamma^2 + 1) / (\kappa + 3)$, where $\gamma$ is skewness and $\kappa$ is excess kurtosis. Values above $5/9$ suggest bimodal distributions.
\end{itemize}

\subsection{Sentiment and Emotion Analysis}

We analyze a stratified sample of 20,000 comments using the \texttt{robertuito} models for Spanish \cite{perez2022robertuito}: a sentiment classifier (positive/neutral/negative) and an emotion classifier (joy, anger, sadness, fear).
For each comment, we compute a polarity score as $\text{polarity} = P(\text{positive}) - P(\text{negative})$.
Results are aggregated at the article level and compared between bias groups.

\subsection{Network Analysis}

We construct a bipartite graph $G = (U, O, E)$ where $U$ is the set of users, $O$ is the set of media outlets, and edges $e_{uo} \in E$ are weighted by the number of comments user $u$ posted on articles from outlet $o$.

\paragraph{Outlet similarity.}
We compute Jaccard similarity between outlets based on their commenter sets: $J(o_i, o_j) = |U_i \cap U_j| / |U_i \cup U_j|$.
High similarity indicates shared audiences.

\paragraph{Community detection.}
We project the bipartite graph onto the outlet dimension and apply the Louvain algorithm \cite{blondel2008fast} to detect communities of outlets connected by shared commenter bases.

\paragraph{User polarization.}
For each user, we compute a bias exposure rate as the fraction of commented articles labeled as biased.
Users are categorized into low ($<0.4$), mixed ($0.4$--$0.6$), high ($0.6$--$0.8$), and very high ($>0.8$) bias exposure groups.
We analyze outlet diversity (number of distinct outlets commented on) across these groups.

\subsection{Predictive Modeling}

To quantify the predictive power of interaction features, we train three classifiers on the 19-dimensional feature vector:
\begin{itemize}
    \item Logistic Regression (interpretable baseline)
    \item Random Forest (non-linear, feature importance)
    \item Gradient Boosting (ensemble, typically best performance)
\end{itemize}

All models use 5-fold stratified cross-validation.
Features are standardized using z-score normalization.
We report accuracy, macro-F1, and ROC AUC, comparing against a majority-class baseline (61.5\% accuracy).

\subsection{Intra-Article Dynamics}

For articles with at least 10 comments, we analyze temporal patterns within comment threads:
\begin{itemize}
    \item \textbf{Early vs.\ late karma}: Mean karma of the first vs.\ last quartile of comments.
    \item \textbf{Karma drift}: The difference between late and early mean karma (negative drift indicates degradation).
    \item \textbf{Karma slope}: Linear regression coefficient of karma on comment position.
    \item \textbf{Reaction speed}: Time (hours) between the first and tenth comment.
\end{itemize}
